\documentclass[11pt,a4paper]{article}
\usepackage[utf8]{inputenc}
\usepackage[T1]{fontenc}
\usepackage{amsfonts}
\usepackage[french]{babel}
\frenchbsetup{StandardLists=true}
\usepackage{enumitem}
\usepackage{amssymb}
\usepackage{amsmath}
\usepackage[left=2cm,right=2cm,top=2cm,bottom=2cm]{geometry}
\usepackage{multirow}
\usepackage[dvipsnames]{xcolor}

\usepackage[T1]{fontenc}
\usepackage{fouriernc}
\usepackage{graphicx}

\usepackage{setspace}
\singlespacing

\usepackage{pstricks}
\usepackage{qrcode}
\usepackage{xkeyval}
\usepackage{pst-labo}

\usepackage{tcolorbox}
\usepackage{multicol}
\usepackage{caption}
\usepackage{fancyhdr}
\pagestyle{fancy}
\renewcommand\headrulewidth{1pt}
\fancyhead[L]{Lycée Denis Diderot }
\fancyhead[R]{Classe de 3PM}
\setcounter{page}{1}\fancyfoot[C]{page \thepage}
\fancyfoot[R]{Année 2025-2026}
\fancyfoot[L]{Alexis RUIZ}
\usepackage{multirow,tabularx,booktabs}
\newcommand{\cadreblanc}[1]{%
\medskip\framebox[\linewidth][c]{%
\rule{0mm}{#1mm}}\par
\medskip}

\usepackage{awesomebox}
\setlength{\aweboxvskip}{0mm}
\usepackage{pifont}
\usepackage{chemmacros}
\chemsetup{modules=all}
\setlength{\columnseprule}{.25mm}
\setlength{\parindent}{0pt}

\renewcommand{\arraystretch}{2}
\newcommand*\acocher{\quitvmode{\fboxsep0pt \fboxrule0.8pt \fbox{\vrule height1.8ex depth.1ex width0pt \vrule height0pt depth0pt width1.9ex }}}

\frenchbsetup{StandardLists=true}

\begin{document}

\begin{center}
 \LARGE{Chapitre 4 - Les ions}
\end{center}

\normalsize

\hrule\
\\[0,2cm]

L'eau que nous buvons, le sel que nous utilisons en cuisine, les batteries de nos téléphones... tous contiennent des ions. Un ion est un atome qui a perdu ou gagné des électrons. Les ions sont présents partout dans notre environnement et jouent un rôle essentiel dans de nombreux phénomènes.

\vspace{3mm}

\textbf{1. Rappels sur l'atome} \\

\textbf{1.1. Structure de l'atome}

\begin{tcolorbox}[colback=red!5!white,colframe=red!75!black]
{
Un atome est constitué d'un noyau au centre, autour duquel des électrons sont en mouvement.

Le noyau contient :
\begin{itemize}
    \item Des protons, portant chacun une charge élémentaire positive (notée +e)
    \item Des neutrons, ne portant pas de charge électrique
    \item Les électrons portent une charge élémentaire négative (notée -e) et tournent autour du noyau.
    \end{itemize}
}
\end{tcolorbox}


\textbf{Propriété : } Un atome est électriquement neutre : il possède autant de protons que d'électrons. \\

\textbf{Exemple : } L'atome de sodium (symbole Na) possède 11 protons dans son noyau, il contient donc aussi 11 électrons. L'atome de chlore (symbole Cl) possède 17 protons et 17 électrons. \\

\textbf{2. Qu'est-ce qu'un ion ?} \\

\textbf{2.1. Définition d'un ion}

\begin{tcolorbox}[colback=red!5!white,colframe=red!75!black]
{
Un \textbf{ion} est un atome (ou un groupe d'atomes) qui a perdu ou gagné un ou plusieurs électrons.

Un ion n'est donc plus électriquement neutre : il porte une charge électrique positive ou négative.
}
\end{tcolorbox}



\textbf{Remarque : } Lors de la formation d'un ion, seul le nombre d'électrons change. Le noyau (donc le nombre de protons) reste identique. \\


\textbf{2.2. Les cations : ions positifs}

\begin{tcolorbox}[colback=red!5!white,colframe=red!75!black]
{
Un \textbf{cation} est un ion qui a perdu un ou plusieurs électrons.

Il possède donc plus de protons que d'électrons : sa charge globale est positive.

On note un cation en écrivant le symbole de l'atome suivi d'un exposant indiquant le nombre de charges positives.
}
\end{tcolorbox}


\begin{tcolorbox}[colback=red!5!white,colframe=red!75!black]
{

\begin{multicols}{2}

\textbf{Exemple : } L'ion sodium \ch{Na+}

\begin{itemize}
    \item L'atome de sodium possède 11 protons et 11 électrons
    \item Il perd 1 électron pour former l'ion sodium
    \item L'ion sodium possède donc 11 protons et 10 électrons
    \item Sa charge est positive : \ch{Na+}
\end{itemize}

\textbf{Exemple : } L'ion cuivre (II) \ch{Cu^{2+}}

\begin{itemize}
    \item L'atome de cuivre possède 29 protons et 29 électrons
    \item Il perd 2 électrons pour former l'ion cuivre (II)
    \item L'ion cuivre (II) possède donc 29 protons et 27 électrons
    \item Sa charge est positive : \ch{Cu^{2+}}
\end{itemize}
}
\end{tcolorbox}

\textbf{2.3. Les anions : ions négatifs}

\begin{tcolorbox}[colback=red!5!white,colframe=red!75!black]
{
Un \textbf{anion} est un ion qui a gagné un ou plusieurs électrons.

Il possède donc plus d'électrons que de protons : sa charge globale est négative.

On note un anion en écrivant le symbole de l'atome suivi d'un exposant indiquant le nombre de charges négatives.
    
\end{multicols}}
\end{tcolorbox}



\begin{tcolorbox}[colback=red!5!white,colframe=red!75!black]
{
\textbf{Exemple : } L'ion chlorure \ch{Cl-}

\begin{itemize}
    \item L'atome de chlore possède 17 protons et 17 électrons
    \item Il gagne 1 électron pour former l'ion chlorure
    \item L'ion chlorure possède donc 17 protons et 18 électrons
    \item Sa charge est négative : \ch{Cl-}
\end{itemize}
}
\end{tcolorbox}



\textbf{3. Les principaux ions à connaître} \\
\begin{multicols}{2}
\textbf{3.1. Les ions monoatomiques}

Les ions monoatomiques sont formés d'un seul atome. Voici les principaux ions monoatomiques à connaître :

\vspace{5mm}

\begin{center}
\begin{tabular}{|l|l|l|}
\hline
\textbf{Nom de l'ion} & \textbf{Formule} & \textbf{Type} \\
\hline
Ion sodium & \ch{Na+} & Cation \\
\hline
Ion potassium & \ch{K+} & Cation \\
\hline
Ion cuivre (II) & \ch{Cu^{2+}} & Cation \\
\hline
Ion fer (II) & \ch{Fe^{2+}} & Cation \\
\hline
Ion fer (III) & \ch{Fe^{3+}} & Cation \\
\hline
Ion zinc & \ch{Zn^{2+}} & Cation \\
\hline
Ion argent & \ch{Ag+} & Cation \\
\hline
Ion chlorure & \ch{Cl-} & Anion \\
\hline
\end{tabular}
\end{center}

\vspace{5mm}

\textbf{3.2. Les ions polyatomiques}

Certains ions sont formés de plusieurs atomes liés entre eux :

\vspace{5mm}

\begin{center}
\begin{tabular}{|l|l|}
\hline
\textbf{Nom de l'ion} & \textbf{Formule} \\
\hline
Ion hydroxyde & \ch{HO-} \\
\hline
Ion sulfate & \ch{SO4^{2-}} \\
\hline
Ion carbonate & \ch{CO3^{2-}} \\
\hline
Ion nitrate & \ch{NO3-} \\
\hline
\end{tabular}
\end{center}

\vspace{5mm}

\textbf{Remarque : } L'ion hydroxyde \ch{HO-} est aussi noté \ch{OH-}. \\

\end{multicols}

\textbf{4. Comment reconnaître certains ions ?} \\

\textbf{4.1. Principe des tests de reconnaissance}

Il existe des tests chimiques qui permettent d'identifier la présence de certains ions dans une solution. Ces tests consistent à ajouter un réactif qui va réagir avec l'ion recherché et former un précipité coloré. \\

\begin{tcolorbox}[colback=red!5!white,colframe=red!75!black]
{
Un \textbf{précipité} est un solide qui se forme dans une solution lorsqu'on ajoute un réactif.

La couleur du précipité permet d'identifier l'ion présent dans la solution.
}
\end{tcolorbox}


\textbf{4.2. Test de l'ion chlorure \ch{Cl-}}

\begin{tcolorbox}[colback=red!5!white,colframe=red!75!black]
{
\textbf{Réactif utilisé : } Solution de nitrate d'argent (\ch{Ag+} + \ch{NO3-})

\textbf{Observation : } Formation d'un précipité blanc qui noircit à la lumière

\textbf{Conclusion : } La solution testée contient des ions chlorure \ch{Cl-}
}
\end{tcolorbox}


\textbf{4.3. Test de l'ion cuivre (II) \ch{Cu^{2+}}}

\begin{tcolorbox}[colback=red!5!white,colframe=red!75!black]
{
\textbf{Réactif utilisé : } Solution d'hydroxyde de sodium (\ch{Na+} + \ch{HO-})

\textbf{Observation : } Formation d'un précipité bleu

\textbf{Conclusion : } La solution testée contient des ions cuivre (II) \ch{Cu^{2+}}
}
\end{tcolorbox}


\textbf{4.4. Test de l'ion fer (II) \ch{Fe^{2+}}}

\begin{tcolorbox}[colback=red!5!white,colframe=red!75!black]
{
\textbf{Réactif utilisé : } Solution d'hydroxyde de sodium (\ch{Na+} + \ch{HO-})

\textbf{Observation : } Formation d'un précipité vert

\textbf{Conclusion : } La solution testée contient des ions fer (II) \ch{Fe^{2+}}
}
\end{tcolorbox}

\textbf{4.5. Test de l'ion fer (III) \ch{Fe^{3+}}}

\begin{tcolorbox}[colback=red!5!white,colframe=red!75!black]
{
\textbf{Réactif utilisé : } Solution d'hydroxyde de sodium (\ch{Na+} + \ch{HO-})

\textbf{Observation : } Formation d'un précipité rouille (brun-orangé)

\textbf{Conclusion : } La solution testée contient des ions fer (III) \ch{Fe^{3+}}
}
\end{tcolorbox}

\textbf{4.6. Test de l'ion zinc \ch{Zn^{2+}}}

\begin{tcolorbox}[colback=red!5!white,colframe=red!75!black]
{
\textbf{Réactif utilisé : } Solution d'hydroxyde de sodium (\ch{Na+} + \ch{HO-})

\textbf{Observation : } Formation d'un précipité blanc

\textbf{Conclusion : } La solution testée contient des ions zinc \ch{Zn^{2+}}
}
\end{tcolorbox}

\textbf{4.7. Tableau récapitulatif des tests}

\vspace{5mm}

\begin{center}
\begin{tabular}{|l|l|l|}
\hline
\textbf{Ion recherché} & \textbf{Réactif} & \textbf{Couleur du précipité} \\
\hline
\ch{Cl-} & Nitrate d'argent & Blanc (noircit à la lumière) \\
\hline
\ch{Cu^{2+}} & Hydroxyde de sodium & Bleu \\
\hline
\ch{Fe^{2+}} & Hydroxyde de sodium & Vert \\
\hline
\ch{Fe^{3+}} & Hydroxyde de sodium & Rouille (brun-orangé) \\
\hline
\ch{Zn^{2+}} & Hydroxyde de sodium & Blanc \\
\hline
\end{tabular}
\end{center}

\vspace{5mm}

\textbf{Remarque : } Attention, les ions chlorure et les ions zinc donnent tous deux un précipité blanc, mais avec des réactifs différents (nitrate d'argent pour \ch{Cl-} et hydroxyde de sodium pour \ch{Zn^{2+}}). \\


\textbf{5. Les solutions ioniques} \\

\textbf{5.1. Définition}

\begin{tcolorbox}[colback=red!5!white,colframe=red!75!black]
{
Une \textbf{solution ionique} est une solution qui contient des ions dissous dans l'eau.

Les ions sont en mouvement dans la solution et ne sont plus liés entre eux.
}
\end{tcolorbox}



\begin{tcolorbox}[colback=red!5!white,colframe=red!75!black]
{
\textbf{Exemple : } L'eau salée

Le sel de cuisine (chlorure de sodium) a pour formule chimique \ch{NaCl}.

Lorsqu'on dissout du sel dans l'eau, les ions sodium \ch{Na+} et les ions chlorure \ch{Cl-} se séparent et se dispersent dans l'eau.

L'eau salée est donc une solution ionique contenant des ions \ch{Na+} et des ions \ch{Cl-}.
}
\end{tcolorbox}


\textbf{5.2. Électroneutralité d'une solution ionique}

\begin{tcolorbox}[colback=red!5!white,colframe=red!75!black]
{
Une solution ionique est toujours \textbf{électriquement neutre}.

Cela signifie que la charge totale apportée par les cations (ions positifs) est égale à la charge totale apportée par les anions (ions négatifs).
}
\end{tcolorbox}


\end{document}
